\chapter{AI 使用说明}
\label{app:ai-usage}

\section*{如何使用本附录}

本附录依据仓库中实际存在的 \path{ai_usage/raw_prompts.md} 与 \texttt{ai\_usage/AI使用记录.md} 整理，只提供角色边界、核验状态和索引。完整长 prompt 不复制进主 PDF，而与主讲义并列提交。若本附录与原始日志不一致，应以版本化原始日志为准并记录修订。

\section{AI 在本项目中的角色}

AI 在本项目中承担结构、草稿和审查辅助：帮助拆解任务、收敛课程范围、组织教学顺序、生成 LaTeX 初稿、执行静态扫描、修复已定位的编译问题，并整理错误与一致性清单。AI 输出始终是待核验材料，不拥有研究决策权。

AI 不是本项目的真实数据提供者、实证结果提供者、投资建议提供者、最终技术责任人或人工复核替代者。日志中的“提交”“静态检查”或“代码检查”也不自动等于内容已被作者采用。

\section{哪些工作由 AI 辅助}

现有日志显示，AI 辅助工作主要包括：

\begin{itemize}
  \item 任务拆解与范围收敛：把广泛蓝图映射到冻结的十章主体；
  \item 教学结构设计：为点时数据、股票池、因子、评价、组合与回测建立前向接口；
  \item LaTeX 初稿：撰写章节、图表、代码清单、模板和附录；
  \item 静态审查：检查章节顺序、字段覆盖、环境配对、标签/引用、禁止命令和构建产物；
  \item 编译修复辅助：根据用户提供的 Overleaf 错误定位 tcolorbox、TikZ 和数学转写问题；
  \item 错误清单：把前视、幸存者偏差、样本选择、账本不守恒等风险转为测试与停止线；
  \item 可读性与一致性检查：统一术语、版本字段、交叉引用和阶段边界。
\end{itemize}

这些工作没有产生真实金融数据，也没有验证三个教学因子的实证有效性。

\section{哪些决定必须由人工负责}

作者或指定复核人必须负责：课程范围与最终采用；数据来源、授权和字段语义；历史指数与证券主键；财务、价格、复权、公司行动和成交口径；研究阈值与停止条件；代码与公式正确性；真实数据运行；回测解释；是否满足交付要求；以及任何对外使用或投资相关决定。

人工还必须审查 AI 是否遗漏上下文、误读接口、引入未经定义对象、制造过密版面或把示意写成事实。AI 不能自我签署“人工核验完成”，也不能用模型内部推理替代可公开的验证证据。

\section{提问如何逐步收敛}

现有记录从 P01 的首轮培训蓝图开始，P02 将广泛课程收敛为十章主体、贯穿案例和标准交付结构，P03 再把用户提供的正式蓝图纳入仓库。随后按研究事件顺序推进：P03-1 建立数据时间边界，P04—P06 完成股票池、因子和预处理，P07—P08 处理单因子章节及 LaTeX 阻塞，P08-ROBUSTNESS 至 P11 完成稳健性、合成、组合和含成本回测，P12 回补第 1 章研究假设，P13 补齐附录与提交说明，P14 更新结项 README，P15 从用户提供的历史对话补录 P01 与 P02，P16 提交并推送结项记录修改。

每轮 prompt 越来越明确地冻结允许修改文件、章节职责、禁止事项、验收和 Git 操作。这种收敛减少了章节间职责漂移，但也意味着早期蓝图的广泛生产化范围不能直接覆盖后来的十章冻结结构；冲突均按当前任务和 \path{AGENTS.md} 的优先级处理。

\section{事实、公式、代码和编译核验}

核验分为不同层次，不能相互替代：

\begin{itemize}
  \item \textbf{事实与制度：}应由官方规则、真实数据字典或责任人确认；日志没有证明任何实际指数、费用或交易制度口径。
  \item \textbf{公式与接口：}通过正文交叉核对、匿名手算、边界条件和代码职责检查；仍需作者复核数学与业务语义。
  \item \textbf{代码：}现有记录主要检查静态结构、禁止字段、确定性、AST 或伪代码接口；仓库当前没有真实数据研究 Notebook 或完整研究代码可供端到端验证。
  \item \textbf{LaTeX：}多条记录明确说明本地没有 XeLaTeX/latexmk；虽然用户任务描述当前 PDF 可编译，仓库日志没有可据此改写为“本地已编译通过”的证据。
  \item \textbf{PDF 视觉：}若能生成 PDF，应渲染逐页检查表格、TikZ、中文、页眉页脚与分页；当前日志未记录完成该人工视觉验收。
\end{itemize}

\section{已知限制与未完成核验}

当前 AI 日志存在以下明确状态：

\begin{itemize}
  \item P01 与 P02 已根据用户提供的历史 ChatGPT 对话补录原始 prompt、附件名称和页面可见输出摘要；历史输出文件未在当前仓库单独保存，且尚未与当前蓝图逐字比对；
  \item P08 包含一个编译修复记录和一个稳健性章节记录，两个编号均保留，不合并或重命名；
  \item 既有原始 prompt 按任务追加，未检测到完全相同的重复任务块；不同任务重复强调 Git、安全和编译要求属于真实原文，不删除；
  \item 既有 AI 使用记录均为“待人工确认”或“人工核验尚未完成”，没有证据支持改写为已审核；
  \item 多数记录明确说明本地缺少 TeX，实际 PDF 页数、版面和交叉引用仍需统一审查；
  \item 仓库当前没有真实金融数据、研究 Notebook、源代码实现、配置文件或实证输出，主体公式和伪代码不能被描述为已验证策略。
\end{itemize}

上述限制只做披露。补录记录必须保留来源和未核验边界，不把页面未显示的模型、绝对日期、文件全文、人工签署、数据结果或编译记录补造成事实。

\section{AI 提问索引}

\begingroup
\scriptsize
\setlength{\tabcolsep}{1pt}
\begin{longtable}{@{}p{1.35cm}p{2.05cm}p{2.15cm}p{2.35cm}p{2.25cm}p{1.75cm}p{1.45cm}p{1.85cm}@{}}
\caption{基于仓库真实日志的 AI 使用索引}\label{tab:appE-ai-index}\\
\toprule
\texttt{prompt\_id} & 任务名称 & 提问目的 & 主要输出 & 修改文件 & 验证方式 & 人工核验状态 & 已知限制 \\
\midrule
\endfirsthead
\toprule
\texttt{prompt\_id} & 任务名称 & 提问目的 & 主要输出 & 修改文件 & 验证方式 & 人工核验状态 & 已知限制 \\
\midrule
\endhead
\path{P01-BLUEPRINT-001} & 首轮培训蓝图 & 定义能力、流程、边界和课程框架 & 19 章首轮蓝图与验收框架 & 历史输出文件 & 对话页面核对 & 文件全文待核验 & 原文件未入库 \\
\path{P02-SCOPE-001} & 范围收敛与重构 & 收敛主体并设计贯穿案例 & 十章结构、三因子案例与质量闸门 & 历史输出文件、后续蓝图 & 对话与结构核对 & 文件全文待核验 & 附件原件未入库 \\
\path{P03-BLUEPRINT-001} & 正式蓝图入库 & 保存用户蓝图并建立索引 & 蓝图与文档索引 & \path{docs/}、AI 日志 & 文件与 Git 检查 & 待人工确认 & 蓝图范围较广 \\
\path{P03-1-001} & 第 2 章 & 建立数据与时间边界 & PIT、时序、字段审计与代码示例 & 第 2 章、AI 日志 & 静态检查 & 尚未完成 & 无本地 TeX、无真实数据 \\
\path{P04-001} & 第 3 章 & 冻结股票池与标签边界 & 五类集合、标签与边界测试 & 第 3 章、AI 日志 & 静态检查 & 尚未完成 & 制度与价格代理待确认 \\
\path{P05-001} & 第 4 章 & 实现三类原始因子 & 定义卡、公式、schema 与核对 & 第 4 章、AI 日志 & 静态与人工样例计划 & 尚未完成 & 财务与价格口径待确认 \\
\path{P06-001} & 第 5 章 & 建立截面预处理基线 & MAD、方向、标准化与暴露 & 第 5 章、AI 日志 & 静态检查 & 尚未完成 & 参数与行业口径待确认 \\
\path{P07-001} & 修复与第 6 章 & 修复 TikZ/页眉并建立单因子评价 & 工程修复、五维证据链 & 工程文件、第 2—6 章、AI 日志 & 静态扫描与 Git & 尚未完成 & 无本地 TeX、无真实数据 \\
\path{P08-001} & LaTeX 阻塞修复 & 修复裸标签和数学转写 & tcolorbox、TikZ 与公式修复 & 第 3—6 章、AI 日志 & 环境、标签和引用扫描 & 尚未完成 & PDF 未在本地验证 \\
\path{P08-ROBUSTNESS-001} & 第 7 章 & 建立样本外与反过拟合流程 & 时间切分、参数邻域与实验台账 & 第 7 章、AI 日志 & 静态与结构检查 & 尚未完成 & 样本、阈值待确认 \\
\path{P09-001} & 第 8 章 & 检查冗余并合成因子 & 共同覆盖、条件分组、消融与等权分 & 第 8 章、AI 日志 & 静态、算术与接口检查 & 尚未完成 & 无真实增量结果 \\
\path{P10-001} & 第 9 章 & 从综合分生成约束目标 & 目标权重、约束、现金与降级 & 第 9 章、AI 日志 & 静态、AST 与约束检查 & 尚未完成 & 资金与约束待确认 \\
\path{P11-001} & 第 10 章 & 建立含成本事件回测 & 订单、成交、持仓、现金和收益账本 & 第 10 章、进度、AI 日志 & 静态、账本与接口检查 & 尚未完成 & 执行、费用、基准待确认 \\
\path{P12-001} & 第 1 章 & 把想法转为可证伪任务 & 假设卡、预注册、停止线与研究地图 & 第 1 章、进度、AI 日志 & 结构、引用与 AST 检查 & 尚未完成 & 无本地 TeX、阈值待确认 \\
\path{P13-A} & 附录 A—C & 建立检索、清单和空白模板 & 术语表、99 项清单、14 份模板 & 附录 A—C、AI 日志 & 字段、标签、引用和结构检查 & 尚未完成 & PDF 页数和版面待验 \\
\path{P13-B} & 附录 D—E & 建立生产概览、AI 说明和提交索引 & 生产流程、日志索引与提交说明 & 附录 D—E、进度、提交说明、AI 日志 & 日志比对与静态检查 & 尚未完成 & 无本地 TeX、人工复核待完成 \\
\path{P14-README-001} & 结项 README & 更新项目状态与验收入口 & 状态、编译、边界与结项清单 & README、AI 日志 & 链接与静态检查 & 尚未完成 & PDF 与代码未交付 \\
\path{P15-RESTORE-001} & 补录 P01/P02 & 恢复历史提问与来源 & 两条原始提问、索引和状态修正 & AI 日志、附录 E、提交说明 & 历史页面与静态核对 & 文件全文待核验 & 历史原件未入库 \\
\path{P16-GIT-001} & 提交结项记录 & 提交并推送当前修改 & README、补录与索引提交 & README、AI 日志、附录 E & Git 与静态检查 & 内容验收待完成 & 无本地 TeX \\
\bottomrule
\end{longtable}
\endgroup

表中“静态检查”仅表示检查了文本或代码结构，不表示实证、编译或人工核验通过。P01 与 P02 是根据用户提供的历史对话补录，来源与未核验项记录在两份 Markdown 日志中。P13-B 与 P13-A 来自同一条完整用户任务；完整原文只在 \path{raw_prompts.md} 保存一次，避免在日志和主 PDF 中重复膨胀。

\section{完整原始提问如何提交}

完整原始提问按任务顺序保存在 \path{ai_usage/raw_prompts.md}，AI 产出摘要、修改文件、验证方式、人工核验状态和限制保存在 \texttt{ai\_usage/AI使用记录.md}。两者均应与主 PDF 和源码仓库并列提交，并保持版本一致。

主 PDF 只包含本索引，不复制数百行 prompt。若需要审查具体任务，应使用 \texttt{prompt\_id} 在两份 Markdown 日志中定位；发现记录缺失时应标记“记录缺失”或“待人工确认”，不能凭记忆重写。

\section{最终责任声明}

AI 参与了范围整理、结构设计、LaTeX 草稿、静态审查和错误修复辅助。作者负责范围决定、内容采用、事实与公式核验、代码和数据验证、人工复核、编译与最终提交。

本教学材料不构成投资建议。匿名案例、空白模板和结构示意不代表真实证券、真实数据库或实证结果；未经真实数据运行和独立核验的部分不得被描述为已验证策略。
